\section{Canonical Reconstruction and Publication Method}
\label{sec:formal-method}

The method separates source routing from output authorization. For source
profile \(s\), let \(E_s(n)\) denote evidence derived from the normalized final
suffix, \(M_s(m)\) evidence derived from the supplied media type, and \(B_s(x)\)
the result of a bounded signature and structural-header probe. These functions
may narrow the eligible source profiles but do not establish that the whole
input belongs to any profile. Their conjunction defines
\begin{equation}
  R_p(x,n,m)=
  \{s\in K^{\mathrm{in}}_p:
    E_s(n)\land M_s(m)\land B_s(x)\land C_{p,s}(n,m)\},
  \label{eq:eligible-routes}
\end{equation}
where \(K^{\mathrm{in}}_p\) is the finite enabled source-profile set and
\(C_{p,s}\) contains explicit alias and compatibility rules. Routing proceeds
only when
$|R_p(x,n,m)|=1$. An empty set denotes unsupported or conflicting input, whereas
a set with more than one member denotes ambiguity. Neither case is resolved
through precedence, extension override, or best-effort fallback.

Normalization accepts one portable basename and one syntactically valid media
type. It rejects path separators, dot components, control or bidirectional
characters, forbidden device names, trailing dots or spaces, wildcard types,
malformed types, and dangerous suffix chains. Alias equivalence is registered
by profile. For example, \texttt{audio/mp4},
\texttt{audio/x-m4a}, and \texttt{audio/m4a} can converge on the same M4A
profile. An alias does not broaden byte acceptance or authorize another ISO
base media profile. HTTP content-type semantics likewise do not turn a supplied
declaration into whole-file evidence \cite{rfc9110}.

For the unique source route \(s\), reconstruction is a total typed operation
\begin{equation}
  \mathsf{Rec}_{p,s}(u)\in
  \mathsf{Result}(\mathcal A,\mathsf{PipelineError}),\qquad
  a=(y,n_o,m_o,k_o,\ell,e)\in\mathcal A,
  \label{eq:reconstruction-relation}
\end{equation}
where \(y\) is the candidate byte string, \(n_o\) is its generated portable
name, \(m_o\) is its declared output media type, \(k_o\) is its output profile,
\(\ell\) is the performed reconstruction level, and \(e\) is bounded stage
evidence. Failure returns a typed non-deliverable error. A successful operation
returns \(\mathsf{Ok}(a)\), not a final verdict, so transformation success is
not confused with authorization.

The performed level records how the output was obtained. Structural output is
serialized from a complete parsed source using an allowlisted component grammar
and may retain selected compressed payload. Semantic output is serialized from
a decoded intermediate representation under policy bounds. Let
\(\Omega_p\subseteq
K^{\mathrm{in}}_p\times
\{\NotProcessed,\Structural,\Semantic\}\times K^{\mathrm{out}}_p\)
be the registered source-level-output relation. It permits format-preserving
routes and explicit conversions such as WAV to MP3 or WebM to MP4. The required
level \(q_{p,s}\) is fixed before dispatch. Acceptance requires
\begin{equation}
  \ell\geq q_{p,s}
  \quad\land\quad
  (s,\ell,k_o)\in\Omega_p.
  \label{eq:level-satisfaction}
\end{equation}
A handler cannot lower this requirement based on input convenience. If a
semantic path is unavailable or fails, a strict semantic profile blocks instead
of returning a structural candidate.

The versioned registry \(\Pi_v\) contains authorizing parsers whose code paths
are separate from the constructing handlers. Each output profile \(k\) names a
nonempty configured set
\(\Pi^{\mathrm{auth}}_{p,v,k}\subseteq\Pi_v\). A parser report \(\rho\) records
the initial offset \(\rho.o\), terminal offset \(\rho.t\), observed profile,
component count, canonicality checks, and bounded observations. Define
\begin{equation}
\begin{split}
  \mathsf{CompleteOK}_{p}(\rho,k,y)\equiv{}&
  \rho.\mathsf{ok}\land \rho.o=0\land\rho.t=|y|\\
  &{}\land \rho.\mathsf{profile}=k
  \land \rho.\mathsf{canonical}_{p,k},
\end{split}
\label{eq:complete-output-report}
\end{equation}
where \(\rho.\mathsf{canonical}_{p,k}\) denotes satisfaction of the
policy-indexed component, ordering, cardinality, integrity, cross-reference,
and resource predicates for \(L_{p,k}\). The configured authorizing condition is
\begin{equation}
  \mathsf{Valid}_{p,v,k}(y)\equiv
  \Pi^{\mathrm{auth}}_{p,v,k}\neq\varnothing
  \land
  \bigwedge_{j\in\Pi^{\mathrm{auth}}_{p,v,k}}
  \mathsf{CompleteOK}_{p}(\rho_j(y),k,y).
  \label{eq:authorizing-validation}
\end{equation}
The reference implementation normally selects one registered authorizing
grammar from the declared output type and then checks that its report agrees
with the generated name and byte recognizer. It does not run every parser in
the registry. External validators from different tool lineages are used only
in the empirical protocol and do not enter the analytical acceptance predicate.

The candidate acceptance predicate is
\begin{align}
  \Accept_{p,v}(u,a) \equiv {}&
  R_p(x,n,m)=\{s\}
  \land \mathsf{Rec}_{p,s}(u)=\mathsf{Ok}(a) \nonumber\\
  &\land \mathsf{InputOK}_p(n,m,s)
  \land (s,\ell,k_o)\in\Omega_p \nonumber\\
  &\land \mathsf{NameOutOK}_p(n_o,k_o)
  \land \mathsf{TypeOutOK}_p(m_o,k_o) \nonumber\\
  &\land \mathsf{ByteTypeOK}_p(y,k_o)
  \land \mathsf{LevelOK}_p(s,\ell) \nonumber\\
  &\land \mathsf{BoundsOK}_p(x,y,e) \nonumber\\
  &\land \mathsf{ScanOK}_p(y)
  \land \mathsf{Valid}_{p,v,k_o}(y).
  \label{eq:extended-acceptance}
\end{align}
Every conjunct is necessary. Signature policy is evaluated on the reconstructed
bytes because scanning only the source would not authorize the emitted object.
An optional pre-scan can reject early but cannot replace the rebuilt-output
scan. Output bytes, expansion ratio, dimensions, frames, duration, channels,
component counts, and nesting are checked with overflow-safe arithmetic.

Reconstruction success creates only a candidate. An unavailable required level
cannot become a weaker pass-through path, and bytes remain unauthorized until
complete output validation. \Cref{fig:reconstruct-or-reject} summarizes this
decision.

\begin{figure}[t]
\centering
\resizebox{0.98\textwidth}{!}{%
\begin{tikzpicture}[
  guard/.style={diamond,draw,aspect=2.0,align=center,fill=blue!7,
    font=\scriptsize,inner sep=2pt},
  transform/.style={draw,rounded corners=2pt,align=center,fill=green!8,
    font=\scriptsize,minimum width=32mm,minimum height=10mm,inner sep=2pt},
  candidate/.style={draw,rounded corners=2pt,align=center,fill=yellow!10,
    font=\scriptsize,minimum width=28mm,minimum height=10mm,inner sep=2pt},
  terminal/.style={draw,rounded corners=2pt,align=center,fill=gray!14,
    font=\scriptsize,minimum width=31mm,minimum height=10mm,inner sep=2pt},
  reject/.style={terminal,fill=red!8,minimum width=270mm},
  arrow/.style={-{Latex[length=2mm]},thick},
  fail/.style={-{Latex[length=2mm]},thin,dashed}
]
\node[candidate] (input) at (-3.0,0) {untrusted\\media \(u\)};
\node[guard] (route) at (0.2,0) {one eligible\\profile?};
\node[guard] (level) at (4.2,0) {required level\\supported?};
\node[transform] (structural) at (8.3,1.2)
  {structural path\\parse and allowlisted reserialization};
\node[transform] (semantic) at (8.3,-1.2)
  {semantic path\\bounded decode and fresh encoding};
\node[candidate] (candidate) at (12.6,0) {candidate only\\performed level recorded};
\node[guard] (authorize) at (16.7,0)
  {all complete-output\\gates pass?};
\node[terminal] (publish) at (20.7,0)
  {validated non-overwriting\\publication};
\node[reject] (reject) at (8.8,-3.35)
  {\(\Reject(e)\): no deliverable bytes and no new target};

\draw[arrow] (input) -- (route);
\draw[arrow] (route) -- node[above,font=\scriptsize]{yes} (level);
\draw[arrow] (level) -- (structural);
\draw[arrow] (level) -- (semantic);
\draw[arrow] (structural) -- (candidate);
\draw[arrow] (semantic) -- (candidate);
\draw[arrow] (candidate) -- (authorize);
\draw[arrow] (authorize) -- node[above,font=\scriptsize]{yes} (publish);

\draw[fail] (route.south) -- node[left,font=\scriptsize]{no} (route.south |- reject.north);
\draw[fail] (level.south) -- node[left,font=\scriptsize]{unsupported} (level.south |- reject.north);
\draw[fail] (structural.south) -- (structural.south |- reject.north);
\draw[fail] (semantic.south) -- (semantic.south |- reject.north);
\draw[fail] (candidate.south) -- (candidate.south |- reject.north);
\draw[fail] (authorize.south) -- node[right,font=\scriptsize]{no} (authorize.south |- reject.north);
\draw[fail] (publish.south) -- node[right,font=\scriptsize]{commit fault} (publish.south |- reject.north);
\end{tikzpicture}%
}
\caption{Fail-closed reconstruction decision. Solid edges form the only
publication path. Dashed edges denote disagreement, unsupported reconstruction,
transformation failure, failed or indeterminate authorization, and commit
failure. All converge on one non-deliverable terminal state.}
\label{fig:reconstruct-or-reject}
\end{figure}

The eight analytical invariants are summarized in
\cref{tab:method-invariants}. Each specifies an observable condition that maps
directly to a test or release gate.

\begin{table}[t]
\centering
\caption{Formal invariants and their empirical interpretation.}
\label{tab:method-invariants}
\small
\begin{tabularx}{\textwidth}{L{0.08\textwidth}L{0.31\textwidth}Y}
\toprule
ID & Invariant & Required observation \\
\midrule
F1 & Complete output consumption & Parser starts at zero and terminates at $|y|$ \\
F2 & Component closure & Every emitted component is admitted at its parsed location \\
F3 & Output evidence convergence & Generated name, declared type, byte recognition, and the authorizing report identify one output profile \\
F4 & Semantic source-independence & No source payload is copied or referenced as an output component on a strict semantic route \\
F5 & Second-pass stability & A second pass is exact for exact profiles or satisfies the registered semantic closure for lossy profiles \\
F6 & Bounded execution & Bytes, structure, time, expansion, and child resources remain within policy \\
F7 & Validation separate from construction & The authorizing parser is distinct from the constructing handler \\
F8 & Clean-result publication & Non-clean and error outcomes expose no bytes, create no new target, and leave any pre-existing target unchanged \\
\bottomrule
\end{tabularx}
\end{table}

Policy comparison is defined \emph{syntactically} on the components that enter
post-reconstruction validation, so that the monotonicity result below is a
theorem about those components rather than a restatement of an extensional
inclusion. Let \(w=(s,a,\{\rho_j(y)\})\) be a fixed post-reconstruction witness,
with \(a=(y,n_o,m_o,k_o,\ell,e)\), and let \(\mathsf{PostOK}_{p,v}(w)\) denote
the conjunction of the source-level-output relation and the output-name,
output-type, byte-type, level, bounds, scan, and authorizing-parser predicates
in \cref{eq:extended-acceptance} (that is, \(\Accept_{p,v}\) with the routing
conjunct \(R_p(x,n,m)=\{s\}\) removed).

\begin{definition}[Syntactic policy refinement]
\label{def:policy-refinement}
Write \(p_2\sqsubseteq p_1\) when, for the fixed output profile \(k_o\) and
parser-registry version \(v\), all of the following hold componentwise:
(i) every scalar bound of \(p_2\) is no larger than the corresponding bound of
\(p_1\), i.e.\ \(\mathcal B_{p_2}\le\mathcal B_{p_1}\) coordinatewise;
(ii) the required reconstruction level is no weaker,
\(q_{p_2,s}\ge q_{p_1,s}\), and \(\Omega_{p_2}\subseteq\Omega_{p_1}\);
(iii) the output language is no broader, \(L_{p_2,k_o}\subseteq L_{p_1,k_o}\),
together with the corresponding inclusions for the admissible name and
media-type sets, \(\mathcal N_{p_2,k_o}\subseteq\mathcal N_{p_1,k_o}\) and
\(\mathcal M_{p_2,k_o}\subseteq\mathcal M_{p_1,k_o}\);
(iv) the signature set is no smaller, \(\Sigma_{p_1}\subseteq\Sigma_{p_2}\),
so that \(p_2\) scans for at least the indicators \(p_1\) does and admits no
scan bypass; and
(v) the authorizing-parser set and canonicality predicate are preserved,
\(\Pi^{\mathrm{auth}}_{p_1,v,k_o}=\Pi^{\mathrm{auth}}_{p_2,v,k_o}\) and
\(\rho.\mathsf{canonical}_{p_1,k_o}\Leftrightarrow\rho.\mathsf{canonical}_{p_2,k_o}\),
and the publication contract of \cref{eq:publish-transition} is unchanged.
\end{definition}

Note the deliberate direction of the signature-set inclusion in
\cref{def:policy-refinement}(iv). The scan predicate \(\mathsf{ScanOK}_p(y)\)
holds when no signature in \(\Sigma_p\) matches \(y\); it is therefore
\emph{antimonotone} in the signature set, and a naive ``larger \(\Sigma\)''
ordering would break the result. Requiring \(\Sigma_{p_1}\subseteq\Sigma_{p_2}\)
is what makes \(\mathsf{ScanOK}\) monotone in the same direction as the other
conjuncts.

\begin{lemma}[Componentwise monotonicity]
\label{lem:component-monotone}
If \(p_2\sqsubseteq p_1\) then each conjunct of \(\mathsf{PostOK}\) is preserved
from \(p_2\) to \(p_1\) on the fixed witness \(w\): for every component
predicate \(\pi\) among the source-level-output, output-name, output-type,
byte-type, level, bounds, scan, and authorizing-parser checks,
\(\pi_{p_2}(w)\Rightarrow\pi_{p_1}(w)\).
\end{lemma}

\begin{proof}
Each conjunct is checked against a set or bound that
\cref{def:policy-refinement} relaxes from \(p_2\) to \(p_1\).
\emph{Bounds:} \(\mathsf{BoundsOK}\) asserts that the recorded scalars of
\((y,e)\) lie at or below the policy bounds; by (i) the \(p_1\) bounds dominate
the \(p_2\) bounds, so satisfaction under \(p_2\) implies satisfaction under
\(p_1\). \emph{Level:} \(\mathsf{LevelOK}\) asserts \(\ell\ge q_{p,s}\) and
\((s,\ell,k_o)\in\Omega_p\); by (ii) \(q_{p_1,s}\le q_{p_2,s}\le\ell\) and
\(\Omega_{p_2}\subseteq\Omega_{p_1}\). \emph{Output language, name, and type:}
\(\mathsf{ByteTypeOK}\), \(\mathsf{NameOutOK}\), and \(\mathsf{TypeOutOK}\)
assert membership of \(y\), \(n_o\), and \(m_o\) in the respective admissible
sets; by (iii) each \(p_2\) set is a subset of the corresponding \(p_1\) set, so
membership is preserved. \emph{Scan:} \(\mathsf{ScanOK}_p(y)\) holds when
\(\Sigma_p\) has no match on \(y\); by (iv) \(\Sigma_{p_1}\subseteq\Sigma_{p_2}\),
so the absence of a \(\Sigma_{p_2}\)-match implies the absence of a
\(\Sigma_{p_1}\)-match. \emph{Authorizing parser:} by (v) the authorizing set
and canonicality predicate are identical under \(p_1\) and \(p_2\), and the
reports \(\{\rho_j(y)\}\) are fixed in \(w\), so \(\mathsf{Valid}_{p_2,v,k_o}(y)
\Leftrightarrow\mathsf{Valid}_{p_1,v,k_o}(y)\). Each conjunct is therefore
preserved.
\end{proof}

\begin{theorem}[Post-reconstruction validation monotonicity]
\label{thm:validation-monotonicity}
For a fixed post-reconstruction witness \(w\) and parser-registry version
\(v\),
\begin{equation}
  p_2\sqsubseteq p_1\land\mathsf{PostOK}_{p_2,v}(w)
  \Longrightarrow \mathsf{PostOK}_{p_1,v}(w).
  \label{eq:validation-monotonicity}
\end{equation}
\end{theorem}

\begin{proof}
\(\mathsf{PostOK}_{p,v}(w)\) is the conjunction of the component predicates in
\cref{lem:component-monotone}. Assuming \(p_2\sqsubseteq p_1\) and
\(\mathsf{PostOK}_{p_2,v}(w)\), every \(p_2\)-conjunct holds; by
\cref{lem:component-monotone} each corresponding \(p_1\)-conjunct holds; their
conjunction is \(\mathsf{PostOK}_{p_1,v}(w)\).
\end{proof}

\Cref{thm:validation-monotonicity} is a post-reconstruction validation
property, not an end-to-end routing property, and it holds only with the
witness \(w\) (in particular the candidate \(y\) and the reports
\(\{\rho_j(y)\}\)) held fixed. A policy change can alter the eligible route set,
select a different reconstruction, or produce different bytes. Such cases are
recorded as route or reconstruction changes and do not instantiate
\cref{eq:validation-monotonicity}.

The result constructor enforces the delivery boundary with an option type:
\begin{equation}
  \mathsf{deliverable}(o,u,a)=
  \begin{cases}
    \mathsf{Some}(y), &
      o=\mathsf{Ok}(\Clean,e)\land\Accept_{p,v}(u,a),\\
    \mathsf{None}, & \text{otherwise}.
  \end{cases}
  \label{eq:clean-only-delivery}
\end{equation}
This rule applies to in-process values, interface buffers, and path publication.
Suspicious, blocked, cancelled, timed-out,
malformed, unsupported, and internally failed operations are distinguishable
in evidence while sharing the same no-deliverable postcondition.

Path publication has a non-overwriting abstract contract. The system derives
the final basename and canonical suffix, stages accepted bytes, revalidates the
policy snapshot and destination condition, and invokes a publication primitive
whose platform-specific guarantees are stated separately. The abstract
transition is
\begin{equation}
  (D,\bot)\xrightarrow{
    o=\mathsf{Ok}(\Clean,e)\land\Accept_{p,v}(u,a)}
  (D\cup\{q\mapsto y\},q),
  \label{eq:publish-transition}
\end{equation}
where \(D\) is the pre-state and \(q\notin\mathrm{dom}(D)\). A failed component
publication call neither creates nor replaces \(q\). If \(q\) already exists or
is created by another actor during the attempt, that entry remains unchanged.
The component does not report a staging object as a final path and removes
staging state on abort to the extent guaranteed by the host filesystem.

The delivery guarantee rests on three assumptions that are stated explicitly
because they, not the acceptance predicate, carry the security weight.

\begin{assumption}[Parser soundness]
\label{asm:parser-soundness}
Every parser \(j\in\Pi^{\mathrm{auth}}_{p,v,k_o}\) is sound for
\(L_{p,k_o}\) in the direction of \cref{eq:complete-profile}: whenever its
report satisfies \(\mathsf{CompleteOK}_{p}(\rho_j(y),k_o,y)\), then
\(y\in L_{p,k_o}\).
\end{assumption}

\begin{assumption}[Strict semantic construction]
\label{asm:semantic-construction}
A handler on a strict semantic route serializes its output solely from the
decoded intermediate representation and never copies or references a byte range
of the source \(x\) as an output component.
\end{assumption}

\begin{assumption}[Publication primitive]
\label{asm:publication}
Publication uses an exclusive-create primitive with the non-overwrite semantics
in \cref{eq:publish-transition}. On failure, the primitive creates no target and
replaces none. If \(q\) already exists or is created concurrently, \(q\)
remains unchanged.
\end{assumption}

\begin{theorem}[Conditional canonical delivery]
\label{thm:conditional-delivery}
Under \cref{asm:parser-soundness}, if \(\Accept_{p,v}(u,a)\) holds then
\(y\in L_{p,k_o}\). Under \cref{asm:semantic-construction}, a strict semantic
route copies or references no source payload as an output component. Under
\cref{asm:publication}, a non-clean result creates or replaces no target through
the component API.
\end{theorem}

\begin{proof}
\(\Accept_{p,v}(u,a)\) contains the conjunct \(\mathsf{Valid}_{p,v,k_o}(y)\)
(\cref{eq:authorizing-validation}), which requires
\(\Pi^{\mathrm{auth}}_{p,v,k_o}\neq\varnothing\) and
\(\mathsf{CompleteOK}_{p}(\rho_j(y),k_o,y)\) for every configured \(j\). Taking
any such \(j\), \cref{asm:parser-soundness} gives \(y\in L_{p,k_o}\). The
payload-independence and publication clauses follow directly from
\cref{asm:semantic-construction,asm:publication}, respectively, together with
\cref{eq:clean-only-delivery}, which withholds bytes on every non-clean result.
\end{proof}

\Cref{thm:conditional-delivery} is deliberately shallow given its assumptions:
the acceptance predicate is engineered so that soundness of the hand-written
authorizing parsers is the sole load-bearing premise of canonical membership.
\Cref{asm:parser-soundness} is therefore the central limitation of the
analytical contribution. It is not discharged here; the differential and fuzzing
protocol of \cref{sec:evaluation} reduces, but cannot prove, the risk of an
unsound parser, and a mechanized correctness proof of the registered grammars
remains future work. \Cref{thm:conditional-delivery} permits coincidental byte
equality and does not exclude harmful decoded meaning, and none of its clauses
proves codec correctness or malware absence.
